\section{Supplementary methodological details}
\label{app:methods}

\subsection{Final predictor inventory}
\label{app:predictors}

The final predictor inventory will document the exact variables used to reproduce the Tuku analysis. For each predictor, the inventory will state its source dataset, temporal transformation, lag or averaging period, and availability during the target month. \placeholder{INSERT THE FROZEN PREDICTOR TABLE FROM THE FINAL TUKU RUN}. Internal experiment labels will remain in the run documentation and will not be used in the manuscript.

\subsection{Model fitting and update settings}
\label{app:model_settings}

The reproducibility record will report the initial calibration period, the six-month block boundaries, the predictor scaling procedure, the fitted regression settings, and the rules used when an input observation was unavailable. \placeholder{INSERT VERIFIED MODEL SETTINGS AND SOFTWARE ENVIRONMENT FROM THE FINAL RUN}. These details will reproduce the reported estimates without interrupting the main methodological narrative.

\subsection{Prediction interval calibration}
\label{app:intervals}

Because the independent section models were fitted using Bayesian ridge regression, the prediction intervals were derived analytically from the posterior predictive distribution. Unlike empirical quantile methods that require an accumulated error archive, the Bayesian formulation provided a theoretically justified interval for every point estimate beginning with the first evaluation block. The prior distribution and noise precision parameters were estimated directly from the available calibration data during each model update. The resulting posterior covariance matrix of the coefficients, combined with the estimated noise precision, defined the uncertainty for each subsequent prediction.

\subsection{Reduced-frequency MLCW measurement settings}
\label{app:sparse_measurement_settings}
\label{app:no_update_settings}

The reduced-frequency analysis contained six scenarios defined by measurement intervals of six or twelve months and initial calibration periods of 36, 60, or 96 consecutive months. Each scenario began at the first complete monthly record and continued through the last complete measurement interval ending within the common MLCW, GWL, and cGNSS observation period. Incomplete intervals at the end of the record were excluded.

For depth section $s$ and measurement interval $I$, the endpoint observation was the accumulated MLCW compaction,
\begin{equation}
Y_{s,I}=\sum_{t\in I}y_{s,t},
\label{eq:interval_compaction}
\end{equation}
where $y_{s,t}$ denotes the hidden monthly compaction increment retained for retrospective evaluation. The endpoint observation constrained the update before the next interval, whereas the individual values of $y_{s,t}$ within $I$ remained unavailable to the model. \placeholder{INSERT THE VERIFIED NUMERICAL FORMULATION AND WEIGHTING OF THE ENDPOINT CONSTRAINT FROM THE FROZEN IMPLEMENTATION}.

Monthly errors were calculated as $\widehat{y}_{s,t}-y_{s,t}$. Endpoint error was defined as accumulated estimated compaction minus accumulated observed compaction,
\begin{equation}
e_{s,I}=\sum_{t\in I}\widehat{y}_{s,t}-Y_{s,I}.
\label{eq:endpoint_error}
\end{equation}
Positive errors represented estimates that were more positive than the observed values. Monthly and endpoint MAE, RMSE, and mean signed error were reported by measurement interval, initial calibration period, and depth section, with pooled values calculated across all six sections. Endpoint totals were not evaluated with the coefficient of determination. The monthly prediction intervals described in \Cref{subsec:predictive_uncertainty} applied only to the primary delayed-data analysis and were not combined across the reduced-frequency intervals.

\FloatBarrier
